%% ma: L12-B08-0119
%% lop: 12
%% chuong: 2
%% bai: 8
%% dang: 12.8.8
%% loai: TN4
%% muc: TH
%% dap_an: "A"
%% nguon: "CĐ2.3.pdf giáo viên gửi 27/10/2026 (M07, Lê Minh Khôi – Nguyễn Anh Khôi, Chương II Chủ đề 3 Biểu thức tọa độ của các phép toán vectơ), Câu 31"
%% kien_thuc: "Bài 8 (điểm chia đoạn thẳng theo tỉ số)"
%% trang_thai: chinh-thuc
\begin{ex}
Trong không gian $Oxyz$, cho hai điểm $A(3;1;-2)$, $B(2;-3;5)$. Tọa độ của điểm $M$ thuộc đoạn $AB$ sao cho $MA=2MB$ là
\choice{\True $M\left(\dfrac73;-\dfrac53;\dfrac83\right)$}{$M(4;5;-9)$}{$M\left(\dfrac32;-5;\dfrac{17}2\right)$}{$M(1;-7;12)$}
\loigiai{$M$ thuộc đoạn $AB$ nên $\overrightarrow{MA}=-2\overrightarrow{MB}\Rightarrow M=\dfrac{A+2B}3=\left(\dfrac73;-\dfrac53;\dfrac83\right)$.}
\end{ex}
